\centerline{\bf \Huge Ура, ТЧ!}

\begin{flushright}
        \scriptsize{}
\end{flushright}

\problem{1}
Операция удвоения цифры натурального числа состоит в умножении этой цифры на $2$ (если это произведение оказывается двузначным, то цифра в следующем разряде числа увеличивается на $1$, как при сложении «в столбик»). Например из числа 9817 удвоениями цифр $7, 1, 8$ и $9$ можно получить числа $9824, 9827, 10617 $и $18817$ соответственно. Можно ли из числа $22 \ldots 22$ ($20$ двоек) несколькими такими операциями получить число $22 \ldots 22$ ($21$ двойка)?

\problem{2}
Взяли пять натуральных чисел и для каждых двух записали их сумму. Могло ли оказаться, что все 10 получившихся сумм оканчиваются разными цифрами?

\problem{3}
Найдите все четные натуральные числа $n$, у которых число делителей (включая 1 и само $n$ ) равно $\frac{n}{2}$. (Например, число 12 имеет 6 делителей: $1,2,3,4,6,12$.

\problem{4}
Найдите все пары простых чисел $p$ и $q$, обладающие следующим свойством: $7p + 1$ делится на $q$, а $7q + 1$ делится на $p.$

\problem{5}
Даны натуральные числа от 1 до 2040, причем $n$ чисел покрашены в зеленый цвет. При каком наибольшем $n$ может оказаться так, что ни одна степень двойки не покрашена и не представима в виде суммы двух зеленых чисел?

\problem{6}
В числе $A$ цифры идут в возрастающем порядке (слева направо). Чему равна сумма цифр числа $9\cdot A$?

\problem{7} Может ли сумма четырёх различных чисел равняться сумме всевозможных попарных НОДов этих чисел?

\problem{8}
На доске написаны $n$ различных целых чисел, любые два из них отличаются хотя бы на 10. Сумма квадратов трёх наибольших из них меньше трёх миллионов. Сумма квадратов трёх наименьших из них также меньше трёх миллионов. При каком наибольшем $n$ это возможно?

\newpage

\hspace{3mm}

\centerline{\bf \Huge О нет, ТЧ...}

\problem{1}
Найдите наименьшее натуральное число, дающее попарно различные остатки при делении на $2, 3, 4, 5, 6, 7, 8, 9, 10$.

\problem{2}
Докажите, что для любых натуральных $x, y$ верно, что $4xy-1$ является делителем $(4x^2-1)^2$ тогда и только тогда, когда $x = y.$

\problem{3}
Пусть $p$-простое число, большее 3. Докажите, что найдётся натуральное число $y$, меньшее $p / 2$ и такое, что число $p y+1$ невозможно представить в виде произведения двух целых чисел, каждое из которых больше $y$.


\problem{4}
Докажите, что найдётся такое натуральное число $n>10^{2025}$, что сумма всех простых чисел, меньших $n$, взаимно проста с $n$.

\problem{5}
Можно ли раскрасить все натуральные числа в два цвета так, чтобы никакая сумма двух различных одноцветных чисел не являлась степенью двойки?